\section{Model}

\subsection{Demand}

One run simulates a single passing period. Time zero is the nominal end of the outgoing class and
$T=900$ seconds is the nominal start of the incoming class. Each classroom holds an outgoing and an incoming class
independently with probability $p_{occ}$; each class has between 15 and 25 students, uniformly at random.
Leavers depart uniformly over the first two minutes. A fraction $q$ of leavers go on to another class in the
building (assigned to an incoming room with probability proportional to remaining seats, and to an adjacent
floor with probability $a$); the rest leave via the lobby. The seats in incoming rooms that transfers do not fill
are occupied by students who arrive from the lobby with a triangular distribution between ten and one minutes before
class start, peaking three minutes before. Every student is ready at an elevator the moment they leave their room
or enter the lobby.

\subsection{Elevators}

Cars are simulated in one-second steps. The four-car bank serves floors 1--9 and the two-car bank floors 1--14.
Passengers whose trip includes a floor above 9 can use only the two-car bank. Cars follow a collective rule:
continue in the current direction while there are riders to drop or calls ahead; otherwise reverse or go idle.
Idle cars go to the nearest unclaimed call. Crucially, cars do \emph{not} skip a floor when full: a car stops for
any call in its direction of travel and then boards as many waiting people as it has room for. This builds in the
empty stops observed in practice; how often they occur in the model (14.5\% of stops in the baseline) is an
output that we have not compared with observation. A car's capacity for a trip is drawn uniformly from
$\{4,5,6\}$, and with backpack compliance $b$ it is increased by two with probability $b$, independently for each
car load.

Timing parameters are 3 seconds between floors (6 seconds from floor~1 to~2), 8 seconds of fixed door time per
stop, and 1 second per rider boarding or alighting. These are estimates, not measurements. They are the
\emph{parameters} of the model. The rules above, and the demand rules of the previous subsection, are its
\emph{structure}, discussed in Section~\ref{sec:struct}.

\subsection{Behavioral measures and the offset}

\emph{Walking.} A rider whose trip starts on floor 2 or 3 and goes down walks with probability $w_d$ (0.8 in the
main experiments). A rider going up to floor 2 walks with probability 0.5 and to floor~3 with probability 0.2,
reflecting the atrium. Walkers take 20 seconds per flight (40 for the atrium) and do not use the elevators.
\emph{Backpacks.} The compliance parameter $b$ above, set to 1 in the main experiments. \emph{Offset.}
Classes on floors 10--14 start and end $o$ minutes earlier than the grid ($o\in\{0,5,10\}$), preserving their
length. A student moving from a lower floor to an upper floor therefore has $15-o$ minutes between classes.

\subsection{Metrics}

For each run we record the mean and 90th-percentile wait (from when a rider is ready to when the rider boards),
the share of students who reach their next class after it starts (walkers included), the share of stops at
which nobody boards or alights, and the time to clear all queues. All figures average 60 replicates with
different random seeds; we report standard errors of the mean for wait and lateness. Wait statistics cover only
students who ride (walkers wait zero and are excluded), lateness includes walkers, and the clearing time is read on
the registrar's clock, from the nominal end of the outgoing class, for both floor groups. The standard errors
measure replicate-to-replicate variation \emph{within the model}. They say nothing about whether the model is
right.

\subsection{Structural assumptions}
\label{sec:struct}

Table~\ref{tab:struct} lists the assumptions that are built into the simulator's rules rather than exposed as
parameters, with where each lives in the code and which results we judge most exposed to it. The last column is
our reasoning, not a test: none of these assumptions has been varied, so we do not know how large any of the
effects is.

{\small
\begin{longtable}{>{\raggedright\arraybackslash}p{1.45in}>{\raggedright\arraybackslash}p{2.6in}>{\raggedright\arraybackslash}p{1.9in}}
\caption{Structural assumptions embedded in the simulator. None is varied in any experiment.}
\label{tab:struct}\\
\toprule
Assumption & What the simulator does & Results most exposed\\
\midrule
\endfirsthead
\toprule
Assumption & What the simulator does & Results most exposed\\
\midrule
\endhead
\bottomrule
\endfoot
Students never leave the queue &
Every student who is not a designated walker waits for an elevator, however long the line. Nobody takes the stairs
on their own, gives up, or leaves early (\texttt{split\_walkers}). The walking measure is applied to a baseline in
which no one walks. &
Baseline wait and lateness; the size of the walking effect; the benefit of any measure that shortens lines.\\
Compliance is a per-car coin flip &
Backpack compliance $b$ adds exactly two seats to a car load with probability $b$, independently of the car, the
floor, and the riders (\texttt{fresh\_capacity}). Compliance is never clustered by class, floor, or time of day, and the
bonus does not depend on the base capacity or on how many people are waiting. &
The size of the backpack effect; its near-linearity in compliance; the ranking of backpacks over walking.\\
Space is the only constraint on boarding &
Everyone in line boards if there is room; nobody declines a crowded car or waits for the next one. Boarding is
first-in-first-out, and an empty car takes the direction of the first person in line (\texttt{serve\_floor}). &
Wait and the share of empty stops.\\
A stop costs only door time &
Travel is a constant 3~s per floor with no acceleration. Everything else about stopping is folded into a fixed
8~s plus 1~s per rider (\texttt{t\_door}, \texttt{t\_pax}). The door-time sweep varies the value of this cost, not its form. &
The cost of empty stops, hence the value of removing riders from lower floors; the gap between short and long trips.\\
Controller and queues &
Collective rule; one shared queue per floor from which any eligible car draws; the two-car bank serves lower floors on
the same terms as the four-car bank; no zoning, no priority for the upper floors. Cars start the period at random
floors, with no backlog from the previous period (\texttt{World}, \texttt{depart}). &
The finding that floor~10 is the worst floor; the offset results; absolute clearing times.\\
Demand is generated independently and in full &
Outgoing and incoming occupancy of a room are independent draws with the same $p_{occ}$. Every seat of every incoming
class is filled, and nobody is absent. Leavers depart uniformly in the first two minutes, with no early or late
dismissal; non-transferring arrivals come only through the lobby on a triangular schedule; transfers choose a room in
proportion to open seats (\texttt{build\_demand}). &
Queue size and shape; lateness; the effect of the offset.\\
Parameters tuned to the story &
$p_{occ}$ was chosen so that the baseline shows a long line without collapse, and the transfer share $q$ and
adjacent share $a$ are not measured. That the baseline resembles the description is therefore not evidence that the
model is right. &
All absolute levels.\\
Lateness is a hard threshold &
A student is late if they arrive after the nominal class start; there is no grace period and no instructor discretion.
Offset classes keep their length. &
Lateness percentages; the offset trade-off.\\
Every changeover is the same &
One 15-minute period, the same for the 50-minute and 75-minute patterns, and for every block of the day. &
Any statement about a particular time of day.\\
\end{longtable}
}

Two points follow. First, the sensitivity analyses in Section~\ref{sec:occ} vary parameters inside this fixed
structure. They show that the ordering of the measures is robust to occupancy, door time, and adjacent-floor
traffic. They do not show that it would survive different student behavior or a different controller. Second, the
effects are not all of one kind. Some assumptions, such as students never leaving the queue, make the baseline look
worse and could shrink the apparent benefit of every measure if they were relaxed. Others, such as the
coin-flip compliance, affect how much each measure achieves. We have not tested any of this, and we have not
tried to guess its net direction.

